\documentclass{rvcepaper}
\begin{document}

% ---------- HS351TA Entrepreneurship & IPR (Improvement CIE) -- Civil dept. header
\begin{iempaper}
\paperinfo{shownote=0,deptline={Department of Civil Engineering},date={12/01/2026},code={HS351TA},sem={V},exam={Improvement CIE},maxmarks={10 + 50},duration={20 Min + 90 Min},
 title={ENTREPRENEURSHIP \& INTELLECTUAL PROPERTY RIGHTS}}
\iemheadB
\ptitle{Part A}
\begin{cietable}{M}{BT}{CO}
\q{1.}{Define Entrepreneurial Marketing and state one distinguishing feature that differentiates it from traditional marketing.}{2}{1}{3}
\q{2.}{Identify any two bases of market segmentation used in the STP framework.}{2}{1}{3}
\q{3.}{List any two sources of entrepreneurial financing available to start-ups.}{2}{2}{3}
\q{4.}{Define an Executive Summary and state any two key purposes it serves in a Business Plan.}{2}{3}{3}
\q{5.}{Identify the Porter's Generic Strategy adopted by a firm that targets a niche market with customized products, and justify your answer briefly.}{2}{2}{3}
\end{cietable}
\ptitle{Part B}
\begin{cietable}{M}{BT}{CO}
\q{1.}{Analyze the key elements of a comprehensive Business Plan. Explain the role and significance of the Executive Summary, Company Description, and Market Analysis in translating a business idea into a viable entrepreneurial venture.}{10}{2}{2}
\q{2.}{Analyse how the 4Ps of Marketing (Product, Price, Place, Promotion) work together in an entrepreneurial context. Illustrate your answer with one Indian start-up case example, highlighting how integration of the marketing mix creates competitive advantage.}{10}{3}{3}
\q{3.}{Evaluate the role of Market Segmentation, Targeting, and Positioning (STP) in developing an effective entrepreneurial marketing strategy. Explain how STP is integrated with the marketing mix, using a real-world business illustration.}{10}{2}{3}
\q{4.}{Explain the major sources of entrepreneurial financing and compare equity financing and debt financing with respect to control, risk, and long-term sustainability of a start-up. Support your answer with suitable examples.}{10}{3}{4}
\q{5.}{Evaluate the importance of Strategic Planning in business development. Discuss how Vision, Mission, Goals, and Objectives, supported by SWOC Analysis, influence the selection of competitive strategies such as Cost Leadership, Differentiation, and Focus Strategy, and enable sustainable growth through Organic Growth, Mergers \& Acquisitions, or Strategic Alliances.}{10}{3}{3}
\end{cietable}
\end{iempaper}

% ---------- IM352IA Operations Management (Improvement CIE)
\begin{iempaper}
\paperinfo{date={12/01/2026},exam={Improvement CIE},maxmarks={10 + 50},sem={V},prog={UG},duration={30 + 90 Min},
 title={Operations Management},code={IM352IA}}
\iemhead
\begin{cietable}{M}{BT}{CO}
\qpart{Part -- A}
\q{1.}{A \blank[1.2cm] shop deals with customized, low-volume production.}{1}{1}{2}
\q{2.}{In mass production, the layout used is \blank[1.2cm] layout.}{1}{1}{2}
\q{3.}{The maximum output a system can produce is called its \blank[1.2cm].}{1}{1}{2}
\q{4.}{The long-term decision of choosing a geographic site for operations is called \blank[1.2cm].}{1}{2}{3}
\q{5.}{Continuous processes are suitable for \blank[1.2cm] production.}{1}{2}{3}
\q{6.}{A bottleneck is a stage that limits the \blank[1.2cm] of the system.}{1}{2}{3}
\q{7.}{The Theory of \blank[1.2cm] focuses on improving the weakest link in a process.}{1}{2}{3}
\q{8.}{MRP stands for \blank[1.2cm] Requirements Planning.}{1}{1}{3}
\q{9.}{The detailed schedule of what is to be produced is called the \blank[1.2cm].}{1}{2}{3}
\q{10.}{ERP stands for \blank[1.2cm] Resource Planning.}{1}{1}{4}
\qpart{Part -- B}
\q{1}{Explain the concepts of economies and diseconomies of scale with appropriate examples.}{10}{2}{2}
\q{2}{List and explain the three major inputs for a Material Requirement Planning system.}{10}{2}{2}
\q{3}{Explain process reengineering and process improvement, highlighting their differences with suitable examples.}{10}{4}{3}
\q{4}{Describe a systematic approach to process analysis with the help of a block diagram.}{10}{2}{2}
\q{5}{A copy center in an office building prepares bound reports for two clients. The center makes multiple copies (the lot size) of each report. The processing time to run, collate, and bind each copy depends on, among other factors, the number of pages. The center operates 250 days per year, with one 8-hour shift. Management believes that a capacity cushion of 15 percent (beyond the allowance built into time standards) is best. It currently has three copy machines. Based on the following table of information, determine how many machines are needed at the copy center.
\qtab{|l|c|c|}{\hline \textbf{Item} & \textbf{Client X} & \textbf{Client Y}\\\hline Annual demand forecast (copies) & 2,000 & 6,000\\\hline Standard processing time (hour/copy) & 0.5 & 0.7\\\hline Average lot size (copies per report) & 20 & 30\\\hline Standard setup time (hours) & 0.25 & 0.40\\\hline}}{10}{3}{3}
\end{cietable}
\end{iempaper}

% ---------- IM353IA Quality Assurance (Improvement CIE)
\begin{iempaper}
\paperinfo{date={13\textsuperscript{th} Jan 2026},exam={IMPROVEMENT CIE},maxmarks={10 + 50},sem={V},prog={UG},duration={30 + 90 Min},
 title={Quality Assurance},code={IM353IA}}
\iemhead
\begin{cietable}{M}{BT}{CO}
\qpart{Part -- A}
\q{1.}{State any two principles of experimentation.}{2}{3}{01}
\q{2.}{What is a factorial experiment?}{2}{2}{01}
\q{3.}{What is MTBF?}{2}{2}{02}
\q{4.}{What are the three common types of failures?}{2}{3}{02}
\q{5.}{What is a parallel system in reliability?}{2}{2}{02}
\qpart{Part -- B}
\q{1}{Consider the following series-parallel configuration:
\par\vspace{4pt}\centerline{\begin{tikzpicture}[x=1cm,y=1cm,font=\small,every node/.style={draw,minimum width=0.45cm,minimum height=0.45cm,inner sep=1pt}]
\node (a) at (1,1.4) {R}; \node (b) at (3.0,2.1) {R}; \node (c) at (3.0,0.7) {R}; \node (d) at (4.7,1.4) {R};
\node (e) at (1,-0.7) {R}; \node (f) at (3.0,-0.7) {R}; \node (g) at (4.7,-0.7) {R}; \node (h) at (7.2,0.35) {R};
\draw[-{Stealth}] (-0.4,0.35) -- (0,0.35); \draw (0,-0.7) -- (0,1.4);
\draw (0,1.4) -- (a) -- (2,1.4); \draw (2,0.7) -- (2,2.1); \draw (2,2.1) -- (b) -- (4,2.1); \draw (2,0.7) -- (c) -- (4,0.7); \draw (4,0.7) -- (4,2.1); \draw (2,1.4) -- (4,1.4) -- (d) -- (5.7,1.4);
\draw (0,-0.7) -- (e) -- (f) -- (g) -- (5.7,-0.7); \draw (5.7,-0.7) -- (5.7,1.4); \draw[-{Stealth}] (5.7,0.35) -- (h) -- (8.3,0.35);
\end{tikzpicture}}\par\vspace{4pt}
\begin{romi}
\item[i)] Derive the expression for system reliability in terms of the component reliabilities, stating clearly the assumption made.
\item[ii)] Compute the System reliability for $R=0.9$ and $R=0.95$.
\end{romi}}{10}{2}{3}
\q{2}{Explain MTBF, MTTF, MTTR relationships with examples.}{10}{3}{3}
\q{3}{Derive and explain reliability for two components in parallel.}{10}{4}{3}
\q{4}{Explain the principles of experimentation in detail.}{10}{4}{4}
\q{5}{Explain six examples of designed experiments in process improvement.}{10}{4}{4}
\end{cietable}
\end{iempaper}

% ---------- IM254TA Financial Accounting & Costing (Improvement CIE)
\begin{iempaper}
\paperinfo{shownote=0,date={14/01/2026},exam={Improvement CIE},maxmarks={10 + 50},sem={V},prog={UG},duration={30 + 90 Min},
 title={Financial Accounting \& Costing},code={IM254TA}}
\iemhead
{\small Note: Answer all the Questions.}\par\vspace{3pt}
\begin{cietable}{M}{BT}{CO}
\qpart{PART - A}
\q{1}{State any two objectives of standard costing.}{2}{1}{3}
\q{2}{List the components of standard cost.}{2}{3}{4}
\q{3}{Define material cost variance. Write its formula.}{2}{3}{4}
\q{4}{Comment on Material Price Variance with formula.}{2}{3}{4}
\q{5}{List the main components of a cost sheet.}{2}{2}{3}
\qpart{PART - B}
\q{1}{Calculate prime cost from the following information:\newline
Opening stock of raw material = Rs. 12,500\newline
Purchased raw material = Rs. 75,000\newline
Expenses incurred on raw material =Rs. 5,000\newline
Closing stock of raw material = Rs. 22,500\newline
Wages Rs. 47,600\newline
Direct expenses Rs. 23,400}{10}{2}{3}
\q{2}{Super Specialities Ltd has to supply 2,40,000 units of a component ``Wizzo'' annually for its valued customer Goenka. The Company has been producing 20,000 units per month by having 12 runs per annum. Its new Finance Manager Arun says that the production should be brought down to 5,000 units per batch and 48 batches should be run per annum. You are required to advise the Company on the following issues given that the Set up Cost per batch is Rs. 75 per batch and Carrying Cost per unit is Rs. 1 per annum.\newline
1. What is the Economic Batch Size?\newline
2. Should the Company continue producing 20,000 units per batch or should it adopt Arun's suggestion?\newline
3. For least cost, how many batches should be run in a year?\newline
4. What should be the interval between each batch?}{10}{4}{3}
\q{3}{In a manufacturing unit, raw material passes through four processes, I, II, III, and IV and the output of each process is the input for the subsequent process. The losses in the four processes are respectively 25\%, 20\%, 20\% and 16 2/3 \% respectively for I, II, III and IV processes of the input. If the end product at the end of the IV process is 40,000 kg, what is the quantity of raw material required to be fed at the beginning of Process I and the cost of the same at Rs. 5 per kg?}{10}{4}{4}
\q{4}{The standard cost card shows the following details relating to the materials needed to \ldots}{10}{2}{4}
\end{cietable}
\end{iempaper}

% ---------- IM355TBB Enterprise Information Systems (Improvement CIE)
\begin{iempaper}
\paperinfo{date={14/01/2026},exam={Improvement CIE},maxmarks={10 + 50},sem={V},prog={UG},duration={30 + 90 Min},
 title={Enterprise Information Systems},code={IM355TBB},note={Answer all the questions}}
\iemhead
\begin{cietable}{M}{BT}{CO}
\qpart{Part -- A}
\q{1.}{List the three main streams that constitute a Supply Chain Management flow.}{1}{1}{CO1}
\q{2.}{State how an EIS system provides a basic operational prerequisite for global supply chain operations.}{1}{2}{CO1}
\q{3.}{Give two examples of common cybersecurity threats specific to an EIS environment.}{2}{2}{CO4}
\q{4.}{Define Agile Project Management based on its delivery characteristics.}{2}{2}{CO4}
\q{5.}{List two primary operational advantages a business gains by adopting a Cloud-based EIS (SaaS model).}{2}{2}{CO4}
\qpart{Part -- B}
\q{1}{Define the terms ``Web-Enabled ERP'' and the ``Extended Enterprise.'' Explain why the transition to web technology was essential for maximizing the value of ERP in a supply chain context.}{10}{2}{CO1}
\q{2}{Clearly state the primary scope and focus of a standalone Enterprise Resource Planning (ERP) system versus a dedicated Supply Chain Management (SCM) system.}{10}{2}{CO2}
\q{3}{Define Artificial Intelligence (AI) from an Enterprise Information Systems (EIS) perspective. List and briefly explain four distinct business functions within an EIS (e.g., SCM, CRM) where AI is being successfully deployed.}{10}{1}{CO4}
\q{4}{What is the Internet of Things (IoT)? And explain its foundational role in modernizing Enterprise Resource Planning (ERP) systems. Detail how IoT enhances real-time visibility in production and supply chain monitoring.}{10}{1}{CO4}
\q{5}{Discuss the concept of ``Cloud-Native'' EIS architecture. Explain why this approach, utilizing micro-services and APIs, is considered more flexible and scalable than traditional monolithic cloud migration methods.}{10}{2}{CO4}
\end{cietable}
\end{iempaper}

\end{document}
